\section{A Bernoulli model and transfer of state moments}
\label{sec:transfer}

Throughout this Section, $A\subset M$ is a maximal abelian subalgebra,
$M$ has separable predual, and $\varphi$ is a faithful normal state with
$A\subset M_\varphi$. We retain the notation of Section~\ref{sec:kernels}:
$N$ is the algebra generated by the groupoid normalizers of $A$,
$E_N:M\to N$ is the $\varphi$-preserving conditional expectation,
$M^\circ=\ker E_N$, and $N_{\rm alg}$ is the algebra of finite sums of
base functions times kernel groupoid normalizers. The Cartan inclusion
$A\subset N$ is represented by a countable nonsingular measured
equivalence relation $\mathcal R$ on $(X,\mu)$, with its scalar cocycle.

The coloring constructed in Section~\ref{sec:coloring} belongs to the
Bernoulli extension of this relation. We first place that extension and
$M$ in a common von Neumann algebra. The resulting model gives a norm
bound for the part of an operator in $M^\circ$. We then show that every
fixed scalar moment involving finitely many model cylinder functions is
a limit of averaged moments in $M$. This is the step that will bring the
model color projections back into $A$.

\subsection{The model and a free compression bound}

Let $\widetilde X$ be the relative Bernoulli extension of $(X,\mu,\mathcal R)$
from Section~\ref{sec:coloring}, with one Lebesgue label at each site of a
relation class. Write $\widetilde\mu$ for its probability measure and
$B=L^\infty(\widetilde X,\widetilde\mu)$. Re-rooting preserves the label
assignment as a function on the sites. In particular, each extended
relation class projects bijectively onto its base class. Let
$\widetilde N$ be the von Neumann algebra of the extended relation,
with the scalar cocycle of $N$ pulled back to the extension.

\begin{lemma}[The Bernoulli free-product model]
\label{lem:bernoulli-model}
The pullback of base kernel operators identifies $N$ with a von Neumann
subalgebra of $\widetilde N$. There is a faithful normal conditional
expectation
\[
 \widetilde E_N:\widetilde N\longrightarrow N
\]
whose restriction to $B$ integrates labels conditional on the base
point. If $\widetilde\varphi$ denotes integration of the diagonal in
$\widetilde N$, then
\[
 (\varphi|_N)\circ\widetilde E_N=\widetilde\varphi.
\]
Moreover, the reduced amalgamated free product
\[
 \mathcal K=M*_N\widetilde N
\]
has a faithful normal conditional expectation $E_N^{\mathcal K}$ onto
$N$, restricting to $E_N$ and $\widetilde E_N$. The state
\[
 \Phi=(\varphi|_N)\circ E_N^{\mathcal K}
\]
is faithful and normal, restricts to $\varphi$ and
$\widetilde\varphi$ on the two constituent algebras, and satisfies
$B\subset\mathcal K_\Phi$. The two constituent algebras are free over
$N$: the expectation of every alternating product of letters from
$\ker E_N$ and $\ker\widetilde E_N$ is zero.
\end{lemma}

\begin{proof}
Use the right-counting direct-integral representations recalled in
Section~\ref{sec:kernels}. Under the bijection between extended and base
classes, a pulled-back base kernel acts by the same orbit operator on
each fiber. Pullback is a normal faithful representation: it is an
amplification over the conditional probability spaces of labels.
Equivalently, matrix coefficients against vector fields on the extension
disintegrate to normal functionals on the original decomposable algebra.
The pulled-back copy of $N$ lies in $\widetilde N$, since this holds for
the generating kernel normalizers.

The isometry from the original representation space to the extension
space that makes a vector field constant in the labels intertwines these
copies of $N$ and sends the original diagonal state vector to the new
one. Compression by this isometry maps $\widetilde N$ normally into
$N$. To see the range assertion, first consider finite sums of
$B$-functions times lifted base groupoid normalizers. Compression
integrates the coefficient functions over the labels at the relevant
sites and hence has its range in $N$. These finite sums form a weakly
dense $*$-algebra: the lifted partial isomorphisms enumerate the extended
relation, and their cuts by extension-measurable sets, together with
scalar phase functions, give its kernel normalizers. Normality of
compression and weak closure of $N$ give the range assertion for all of
$\widetilde N$.

The compression is unital, $N$-bimodular, and equal to the identity on
$N$, so it gives the asserted normal conditional expectation. It
preserves the diagonal state. That state is faithful by the faithful
normal diagonal expectation for a countable nonsingular relation
algebra. Consequently $\widetilde E_N$ is faithful as well. The diagonal
expectation is $B$-bimodular and $B$ is abelian, so
$B\subset\widetilde N_{\widetilde\varphi}$.

Both expectations onto $N$ are faithful and normal, and the constituent
algebras are $\sigma$-finite. The von Neumann algebra reduced
amalgamated free-product construction therefore gives $\mathcal K$,
$E_N^{\mathcal K}$, generation by the two copies, and their freeness over
$N$; see \cite[Section~2]{Ued99} and \cite[Section~2]{Ued13}. These
hypotheses are the faithful-normal conditional-expectation hypotheses
of that construction, and do not require tracial states. Faithfulness
and normality of $\Phi$ follow from those of the expectation and of
$\varphi|_N$.

For completeness, the product has $\Phi$-preserving conditional
expectations onto each constituent algebra. Fix one constituent and
compress the state representation of $\mathcal K$ to the closure of
its state vectors. The linear span of $N$ and alternating centered
words is weakly dense. Every reduced word containing a centered letter
from the other constituent compresses to zero: test between vectors
from the retained algebra, expand at the two ends, and absorb the
resulting $N$-coefficients into adjacent letters; freeness then applies.
The retained algebra acts normally and faithfully on its state-vector
space. Compression is thus a normal bimodular state-preserving map
onto that represented algebra. By Takesaki's Theorem
\cite{Tak72}, the modular group of $\Phi$ restricts to the modular group
of each constituent. This also agrees with
\cite[Equation~(2.1)]{Ued13}. Since $B$ is in the centralizer of
$\widetilde\varphi$, it is in $\mathcal K_\Phi$.
\end{proof}

\begin{lemma}[Free compression]
\label{lem:free-compression}
In the model of Lemma~\ref{lem:bernoulli-model}, let $p\in B$ be a
projection such that $\widetilde E_N(p)=\lambda1$, where
$0\le\lambda\le1$. For every $w\in M^\circ$,
\begin{equation}
 \|pwp\|\le(\lambda+2\sqrt\lambda)\|w\|.
 \label{eq:free-compression}
\end{equation}
\end{lemma}

\begin{proof}
In $L^2(\mathcal K,\Phi)$, let $P$ be the orthogonal projection onto the
closed span of reduced-word vectors whose first letter is a centered
$M$-letter. The orthogonal complement is spanned densely by the
$N$-vectors and the reduced-word vectors whose first letter is a
centered $\widetilde N$-letter. Both the spanning assertion and the
orthogonality follow from generation and freeness over $N$.

Left multiplication by $w$ sends each of these generators of the
orthogonal complement into the range of $P$. Thus
\[
 (1-P)w(1-P)=0.
\]
On a word beginning with a centered $M$-letter, left multiplication by
$p-\lambda1$ inserts a centered $\widetilde N$-letter at the beginning.
Its image is orthogonal to the range of $P$. Hence
\[
 PpP=\lambda P,
 \qquad
 \|Pp\|^2=\|PpP\|\le\lambda.
\]
Insert $P+(1-P)$ on both sides of $w$ in $pwp$. The term with
$(1-P)w(1-P)$ vanishes. The remaining three terms have norms at most
$\lambda\|w\|$, $\sqrt\lambda\|w\|$, and
$\sqrt\lambda\|w\|$, respectively. This proves
Equation~\eqref{eq:free-compression}.
\end{proof}

\subsection{Finite symbols and their evaluation in the masa}

We now specify the model functions that can be evaluated inside $A$.
Fix an integer $s\ge2$, and partition each Lebesgue label into $s$
equally likely symbols, identified with the $s$th roots of unity. A
finite dyadic symbol partition, with $s$ a power of two, will suffice
when we later approximate arbitrary labels.

A \emph{finite cylinder function} is a bounded function on
$\widetilde X$ that depends measurably on the base point and on these
symbols at a fixed finite list of translated sites. Each translated
site is specified by a partial isomorphism with graph in $\mathcal R$;
the function may also depend on whether the base point belongs to that
map's domain. Let $\mathcal Y$ be the algebra of finite sums of such
functions times lifted base kernel groupoid normalizers. Covariance
and composition show that this is a $*$-algebra. Conditional label
integration gives
\[
 \widetilde E_N(\mathcal Y)\subset N_{\rm alg}.
\]

Choose increasing finite Borel partitions $\mathcal P_j$ of $X$ that
generate $A$ and separate points outside the fixed null exceptions.
For one such partition write $\mathcal P=\{e_a\}_a$, identifying its
pieces with their projections in $A$. For independent uniform $s$th
roots of unity $z_a$, define
\[
 u_{\mathcal P}=\sum_a z_a e_a\in A.
\]
Replace the symbol at a translated site by the value of
$u_{\mathcal P}$ at that site, and replace every lifted base normalizer
by its original normalizer in $N$. This defines an evaluation map,
denoted by $\mathrm{ev}_{\mathcal P}\colon\mathcal Y\to N_{\rm alg}$.
Finite cylinder functions themselves evaluate to elements of $A$.

This evaluation respects multiplication and adjoints and fixes
$N_{\rm alg}$. Substitution is consistent under re-rooting and therefore
with covariance. It also respects almost-everywhere identities. Indeed,
conditional on almost every base point, each assignment of symbols to
any finite list of distinct sites has positive probability. An identity
of cylinder functions therefore holds for every such assignment outside
a null set of base points, including the assignments produced by
$u_{\mathcal P}$. The same reasoning applies to identities between
finite sums of cylinder coefficients times normalizers: check the
finitely supported root columns in the counting representations,
including the common cocycle phases. Equality of these columns almost
everywhere gives equality of the operators by faithfulness of the
diagonal state. In particular, the definition does not depend on the
chosen finite expression.

The norm of an evaluated cylinder coefficient is bounded by its
original essential bound. For every fixed finite normalizer sum,
its expression consequently gives a uniform bound on all evaluations.
No uniform bound over arbitrary finite expressions will be needed.

Before stating the transfer result, we isolate the uniform approximation
estimate needed for its fixed scalar words. It holds for any faithful normal
state and will also be used in the final clipping argument.

\begin{lemma}[Uniform right multiplication in the state norm]\label{lem:state-multipliers}
Let $\varphi$ be a faithful normal state on a von Neumann algebra $M$.
For every $C<\infty$, every fixed $z\in M$, and every $\eta>0$, there
is $\delta>0$ such that
\begin{equation}\label{eq:uniform-fixed-right}
 \|D\|\leq C,\quad\|D\|_\varphi<\delta
 \quad\Longrightarrow\quad \|Dz\|_\varphi<\eta.
\end{equation}
The choice of $\delta$ is uniform over all such $D$.
For $U,D\in M$ and $b\in M_\varphi$,
\begin{equation}\label{eq:state-left-centralizer}
 \|UD\|_\varphi\leq\|U\|\|D\|_\varphi,
 \qquad
 \|Db\|_\varphi\leq\|b\|\|D\|_\varphi.
\end{equation}
More generally, the right multiplier bound is uniform when $z$ ranges
over a norm-bounded set that is precompact in state norm.
Consequently, a bounded family of errors tending uniformly to zero in
state norm remains so after right multiplication by a fixed finite word
whose factors range over such precompact sets or over bounded subsets
of the centralizer.  In particular, any specified bounded sequence
converging to a fixed operator in state norm can be used as a right
factor uniformly over its entire sequence.  These conclusions permit
bounded strong-* approximation of fixed letters in fixed-length scalar
words, uniformly in the intervening bounded centralizer factors and in
the independently varying indices of other specified approximation
sequences.
\end{lemma}

\begin{proof}
Let $\Omega$ be the GNS state vector. Faithfulness makes it separating
and hence cyclic for $M'$.  For any $\alpha>0$, choose $c'\in M'$ with
$\|z\Omega-c'\Omega\|<\alpha$.  Then every $\|D\|\leq C$ satisfies
\[
 \|Dz\|_\varphi
 \leq C\alpha+\|Dc'\Omega\|
 = C\alpha+\|c'D\Omega\|
 \leq C\alpha+\|c'\|\|D\|_\varphi.
\]
Choosing $\alpha$ first proves the uniform assertion
\eqref{eq:uniform-fixed-right}, including families indexed by additional
parameters.

The left multiplier bound is immediate.  For $b\in M_\varphi$ and
$D_0\geq0$, centrality gives
$\varphi(b^*D_0b)=\varphi(D_0bb^*)$.  The functional
$c\mapsto\varphi(D_0c)$ is positive on $M_\varphi$, since for
$c\geq0$ there it equals $\varphi(c^{1/2}D_0c^{1/2})$.
It follows that
$\varphi(b^*D_0b)\leq\|b\|^2\varphi(D_0)$, proving
Equation~\eqref{eq:state-left-centralizer}.

For a norm-bounded state-norm precompact set $\mathcal C$, take a finite
state-norm approximation $z_1,\ldots,z_m$ at a prescribed accuracy.
If $z\in\mathcal C$ is close to $z_i$, then
\[
 \|Dz\|_\varphi
 \leq\|Dz_i\|_\varphi+C\|z-z_i\|_\varphi.
\]
The finitely many fixed-factor moduli give one modulus for
$\mathcal C$.  The set consisting of a convergent sequence and its
limit is precompact, so this applies to every bounded sequence
$z_j$ with $\|z_j-z\|_\varphi\to0$.
Apply the right multiplier bounds successively in the order in which
the factors occur.  At each step the errors remain uniformly bounded
in operator norm, so the next modulus continues to apply.
For a scalar word with a middle error $D$, the estimate
\[
 |\varphi(UDV)|\leq\|U\|\,\|DV\|_\varphi
\]
and telescoping reduce the approximation to the preceding bounds.
Strong-* approximation ensures state-norm convergence of a replaced
letter and its adjoint whenever both occur.  For each fixed approximant
the finite right words may be expanded as needed; its approximation
index can then be sent to infinity after the uniform error estimates.
\end{proof}

\begin{proposition}[Transfer of fixed scalar moments]
\label{prop:moment-transfer}
Let $A\subset M$, $\varphi$, and $\mathcal K$ be as in
Lemma~\ref{lem:bernoulli-model}. Fix a finite symbol size $s$, the
corresponding cylinder algebra $\mathcal Y$, and the generating
partition sequence $(\mathcal P_j)$ just described. For any fixed
integer $k\ge1$ and fixed factors $Z_1,\ldots,Z_k$, each chosen in
$M$ or in $\mathcal Y$, define $Z_h^{(j)}=Z_h$ for a factor chosen in
$M$ and
$Z_h^{(j)}=\mathrm{ev}_{\mathcal P_j}(Z_h)$ for a factor chosen in
$\mathcal Y$. Use the same independent symbols for all evaluations
in a word. Then
\begin{equation}
 \lim_{j\to\infty}
 \mathbb E_{\{z_a\}}
 \varphi\bigl(Z_1^{(j)}\cdots Z_k^{(j)}\bigr)
   =\Phi(Z_1\cdots Z_k).
 \label{eq:moment-transfer}
\end{equation}
\end{proposition}

\begin{proof}
We first record how operator approximation will be used. We then reduce
to centered cyclic words. Their label averages are sums of terms with
prescribed equality relations among partition indices. A shortest pair
of linked occurrences will give a state Cauchy--Schwarz estimate whose
small factor is controlled by Proposition~\ref{prop:atomless-kernels}.

\paragraph{Uniform approximation in fixed words.}
Lemma~\ref{lem:state-multipliers} controls bounded errors in the state
norm when the right word is formed from fixed operators and bounded
varying $A$-functions. Expanding each fixed finite normalizer sum when
needed, telescoping and
\[
 |\varphi(UDV)|\le\|U\|\,\|DV\|_\varphi
\]
bound scalar word errors uniformly over the evaluated diagonal functions.
If a fixed $N$-letter is approximated by a bounded strongly convergent
sequence from $N_{\rm alg}$, that sequence and its limit form a set
precompact in the state norm, so the lemma applies to the whole sequence
as right factors. Each normalizer sum is expanded only after that approximant
is fixed; no uniform bound on all such expansions is needed. The order is
to fix an approximant, pass to the partition limit, and then let its
approximation error tend to zero. The same statements apply in the model,
with $\Phi$ and bounded factors from $B$, which lies in its centralizer.

\paragraph{Reduction to centered alternating words.}
We induct on the number of factors. A word of length one in $M$ gives
an identity. For one factor in $\mathcal Y$, average its cylinder
coefficients. At almost every base point, distinct sites in a fixed
finite translated list eventually belong to distinct atoms of
$\mathcal P_j$. The averaged coefficient thus converges boundedly
almost everywhere to its conditional label average. This is also
state-norm convergence, and the fixed normalizers are allowed right
multipliers. It proves Equation~\eqref{eq:moment-transfer} in length
one.

Adjacent factors from the same algebra can be folded into one factor.
For an alternating word, subtract and add the $N$-expectation of each
factor in the model. Make this same expansion on the evaluated side;
it is not necessary that a centered model factor remain centered after
evaluation. The expectations of $\mathcal Y$-factors lie in
$N_{\rm alg}$ and are fixed by evaluation. Every term containing an
expectation can be reduced in length by folding it into a neighbor.
If an arbitrary $N$-element from the expectation of an $M$-factor must
be multiplied into a $\mathcal Y$-factor, first approximate it boundedly
in the strong-$*$ topology by $N_{\rm alg}$. Such approximation follows
from generation and Kaplansky density, and is legitimate on both sides
by the uniform word estimates above. The induction hypothesis therefore
handles all terms except fully centered alternating words. Their model
expectation is zero by freeness.

We may approximate each centered $M$-factor boundedly in the strong-$*$
topology by an entire analytic element for $\sigma^\varphi$ belonging
to $M^\circ$. For example, convolve its modular orbit with Gaussian
approximate identities; complex translations of the Gaussian give the
entire continuation. The expectation $E_N$ commutes with the modular
group, so the convolutions remain centered. Their complex modular
translates also remain in $M^\circ$, by analytic continuation after
applying $E_N$.

Each centered $\mathcal Y$-factor is a finite sum of terms $vb$, where
$v$ is a base kernel groupoid normalizer and $b$ is a cylinder function
whose conditional label average is zero. To obtain this form, move each
coefficient to the right of its normalizer by covariance on the relevant
supports, incorporating bounded base coefficients into $b$, and subtract
the expectation of each summand. The sum of the subtracted expectations
is zero. Cut the initial support of each $v$ as in
Section~\ref{sec:kernels} to make it entire analytic. These cuts preserve
centering of the cylinder coefficients and have arbitrarily small
state errors in the fixed words. The cylinder functions and their
evaluations belong to the respective centralizers. Hence complex
modular shifts of $vb$ are compatible with evaluation: the shift only
changes the cut normalizer by the same bounded base formula in both
algebras. We fix these analytic approximants before passing to the
partition limit; their complex-shift norms may enter all subsequent
bounds.

For an odd alternating word of length at least three, the first and
last factors lie in the same constituent algebra. Apply the KMS
identity
\[
 \varphi(ZY)=\varphi\bigl(Y\sigma^\varphi_{-i}(Z)\bigr)
\]
to the first evaluated analytic factor, and the corresponding identity
for $\Phi$ in the model. Fold the last factor with the shifted first
factor and apply the induction hypothesis to the shorter word. Its
model value equals that of the original centered word and is therefore
zero.

For an even alternating word, rotate in the same way if necessary to
start in $M^\circ$. Expand the other-side factors as $vb$ and absorb
each $v$ into the preceding centered $M$-factor. Modular shifts of
normalizers, if present, only add fixed base coefficients. Writing the length of this word as $2m$, it remains
to prove that the averaged evaluations of
\begin{equation}
 \varphi(y_1d_1\cdots y_md_m)
 \label{eq:cyclic-centered-word}
\end{equation}
tend to zero. Here the $y_h$ are fixed bounded entire analytic elements
of $M^\circ$, and each $d_h\in A$ is the substitution in a centered
finite cylinder function. All constants below may depend on this fixed
word and on the fixed analytic approximants.

\paragraph{Fourier expansion and index equalities.}
Split each cylinder domain according to the domains of its translated
sites and the equality pattern among those sites. On each piece, its
distinct label symbols are independent. Fourier expansion on the finite
product of cyclic groups expresses the centered coefficient as a finite
sum of monomials with fixed bounded $A$-coefficients. Centering removes
the constant Fourier term. Delete occurrences having exponent zero
modulo $s$. Each remaining monomial has at least one occurrence, and
all its exponents are nonzero modulo $s$. Within each diagonal block
$d_h$, the sites of its occurrences are pointwise distinct on that
block's indicated support.

We may, after further finite splitting, regard every translation used
at an occurrence as an automorphism of the full abelian algebra $A$.
Here the extension away from the specified support need not belong to
$\mathcal R$. Proposition~\ref{prop:atomless-kernels} shows that the
$y_h$ annihilate the atomic part of $A$ on both sides, so only diffuse
supports contribute to Equation~\eqref{eq:cyclic-centered-word}.
The partial isomorphisms map these supports into the diffuse part
modulo null sets. Split a nonnull common domain piece of a block into
two positive-measure pieces. On either smaller piece, both its domain
and the range of each restricted map have nonnull diffuse complements
in the diffuse part: the other piece and its image provide such
complements. Normalize the measures on each pair of complements
separately and use the isomorphism theorem for atomless standard
probability spaces \cite[Theorem~17.41]{Kec95}. It extends each
restricted map to a measure-class automorphism of the diffuse part.
Use the identity on the atomic part. This gives finitely many terms,
and for an occurrence $o$ we can write its translation as
\[
 \lambda_o(f)=f\circ\gamma_o,\qquad f\in A,
\]
where $\gamma_o$ is a nonsingular measurable bijection modulo null sets.

Expand Equation~\eqref{eq:cyclic-centered-word} in partition indices
$a_o$, one for each occurrence. Averaging the independent symbols
$z_a$ keeps precisely the index equality patterns for which the sum
of the Fourier exponents in each equality class is zero modulo $s$.
Since individual exponents are nonzero, every such class has at least
two occurrences. Inclusion--exclusion over the finitely many possible
coarsenings removes inequalities between distinct classes. Thus it
suffices to prove vanishing for a sum with prescribed equalities within
the classes of a partition of the occurrences, every class having size
at least two, and with no inequalities imposed between classes. An
occurrence contributes the projection $\lambda_o(e_{a_o})$.

We will keep selected pair equalities as explicit summation indices.
Whenever the selected pairs form a forest in each equality class,
extend them to a spanning tree of that class. Encode the other tree edges with
independent auxiliary phases. For an edge between indices $a,b$, its
equality indicator is
\[
 \mathbf1_{\{a=b\}}
   =\mathbb E_{\boldsymbol\zeta}
      \bigl(\zeta_a\overline{\zeta_b}\bigr),
\]
where the $\zeta_a$ are independent uniform circle phases; use a fresh
phase vector for each edge. At an occurrence $o$, let $w_o(a)$ be the
product of its incident auxiliary edge phases, with conjugates according
to their orientations and with empty product one. Thus $|w_o(a)|=1$.
Set
\[
 U_o=\sum_a w_o(a)\lambda_o(e_a).
\]
The projections $\{\lambda_o(e_a)\}_a$ form a partition, so $U_o$ is a
diagonal unitary. Summing an unretained index gives $U_o$, while at a
retained index $j$ one has
\[
 w_o(j)\lambda_o(e_j)=U_o\lambda_o(e_j).
\]
Thus each $U_o$, as an operator, is independent of the retained numerical
indices. Different $U_o$ may be correlated through the auxiliary phases;
the estimates below hold uniformly over their joint realizations as
diagonal unitaries before averaging.
Products of these factors commute within each diagonal block.

\paragraph{A pair in one diagonal block.}
Suppose two occurrences in one equality class belong to the same block.
Retain their pair equality and encode the remaining tree edges as just
described. Summing the pair index produces the projection
\[
 p_{\mathcal P}
  =\mathbf1_{\mathrm{support}}
    \sum_a\lambda_o(e_a)\lambda_{o'}(e_a).
\]
The two sites are distinct on the fixed support, while the partitions
separate points. Hence $\mu(p_{\mathcal P})\to0$. Use KMS rotation to
place this projection first in the word. Only fixed analytic letters
are shifted; diagonal factors are in the centralizer and remain
unchanged. The rotated word has the form $\varphi(p_{\mathcal P}Z)$
with a uniform norm bound on $Z$. Centrality gives
\[
 |\varphi(p_{\mathcal P}Z)|
 =|\varphi(p_{\mathcal P}Zp_{\mathcal P})|
 \le\|Z\|\,\mu(p_{\mathcal P}),
\]
so these sums tend to zero uniformly in the auxiliary phases.

\paragraph{A shortest pair between different blocks.}
It remains to treat equality patterns with at most one occurrence of
any class in each diagonal block. This case has $m\geq2$. Number the
blocks by $\mathbb Z/m\mathbb Z$, and for distinct positions $h,k$ define
\[
 d_m(h,k)=\min\{(k-h)\bmod m,(h-k)\bmod m\},
\]
where both residues are taken in $\{1,\ldots,m-1\}$. Among pairs of
occurrences in the same class, choose one minimizing this cyclic distance,
and call the minimum $\ell$. Choose the endpoint order so that one shortest
arc runs forward in the cyclic order of the word, and label its blocks
$0,1,\ldots,\ell$; either endpoint order is allowed in an antipodal tie.

For each interior block $h=1,\ldots,\ell-1$, choose one occurrence and
one other occurrence of its class. Its mate lies outside the entire chosen
arc, including the endpoints. Indeed, a mate in another block $k$ on that
arc would satisfy $d_m(h,k)\leq|h-k|<\ell$, contradicting minimality.
The same inequality shows that the interior classes are distinct from each
other and from the endpoint class. When $\ell=1$ there are no interior
blocks.

Retain the endpoint pair and these interior-to-mate pairs. These pairs
lie in distinct equality classes, so they form a forest and can be
included in the spanning trees used above. The retained indices are
$J=(j_0,j_1,\ldots,j_{\ell-1})$, with $j_0$ linking the endpoints.
Write
\[
 E_0=\lambda_0(e_{j_0}),\qquad
 E_\ell'=\lambda_\ell(e_{j_0}),\qquad
 E_h=\lambda_h(e_{j_h})\quad(1\le h<\ell).
\]
Each $j_h$ for $h\ge1$ occurs in one further retained projection in an
exterior block. Several such mate projections may lie in the same
exterior block.

Rotate the summands by KMS to begin at $E_0$, and insert $E_0$ at the
far end inside the state, using centrality. Absorb fixed diagonal
coefficients into neighboring fixed letters. More explicitly, within
a diagonal block place the retained projections first, then the
auxiliary unitaries, then its fixed coefficient; absorb that coefficient
into the following intervening letter. The fixed intervening letters,
including any modular translates created by rotation, remain bounded
analytic elements of $M^\circ$. Relabel them as $y_h$, and write $V_h\in A$ for the auxiliary unitaries,
which are independent of the explicit indices. Split the word at
$E_\ell'$ using $(E_\ell')^2=E_\ell'$. Before averaging the auxiliary
unitaries, the resulting summands have the form $\varphi(X_JY_J)$, where
\begin{align}
 X_J
  &=E_0V_0
       \left(\prod_{h=1}^{\ell-1}y_hE_hV_h\right)
          y_\ell E_\ell',
 \notag\\
 Y_J
  &=E_\ell'V_\ell
       B_{j_1,\ldots,j_{\ell-1}}E_0.
 \label{eq:shortest-pair-factors}
\end{align}
The product in $X_J$ is in increasing order and is omitted when
$\ell=1$. In $B_{j_1,\ldots,j_{\ell-1}}$, each exterior mate projection
for these indices occurs exactly once; every other factor is independent
of them. Diagonal factors in the same block commute.

Figure~\ref{fig:shortest-pair} illustrates the selected links in one
allowed cyclic pattern.
\begin{figure}[ht]
\centering
\begin{tikzpicture}[
  x=1cm,y=1cm,
  block/.style={draw,rounded corners=1pt,minimum width=1.1cm,minimum height=0.75cm,inner sep=2pt},
  word/.style={-{Stealth[length=1.5mm]},line width=0.35pt},
  equality/.style={densely dashed,line width=0.7pt},
  every node/.style={font=\small}
]
\node[block] (d0) at (0,1.5) {$d_0$};
\node[block] (d1) at (2,1.5) {$d_1$};
\node[block] (d2) at (4,1.5) {$d_2$};
\node[block] (d3) at (6,1.5) {$d_3$};
\node[block] (d4) at (6,-1) {$d_4$};
\node[block] (d5) at (4,-1) {$d_5$};
\node[block] (d6) at (2,-1) {$d_6$};
\node[block] (d7) at (0,-1) {$d_7$};

\draw[word] (d0.east) -- (d1.west);
\draw[word] (d1.east) -- (d2.west);
\draw[word] (d2.east) -- (d3.west);
\draw[word] (d3.south) -- (d4.north);
\draw[word] (d4.west) -- (d5.east);
\draw[word] (d5.west) -- (d6.east);
\draw[word] (d6.west) -- (d7.east);
\draw[word] (d7.north) -- (d0.south);

\coordinate (o51) at ($(d5.north)+(-0.2,0)$);
\coordinate (o52) at ($(d5.north)+(0.2,0)$);
\draw[equality] (d0.north) to[out=90,in=90]
  node[above,pos=0.5] {$j_0$} (d3.north);
\draw[equality] (d1.south) to[out=-65,in=110]
  node[left,pos=0.45] {$j_1$} (o51);
\draw[equality] (d2.south) to[out=-70,in=70]
  node[right,pos=0.45] {$j_2$} (o52);
\foreach \point in {d0.north,d3.north,d1.south,d2.south,o51,o52}
  \fill (\point) circle[radius=1.4pt];

\node[fill=white,inner sep=2pt] at (3,2.1) {short arc: $X_J$};
\node at (3,-1.7) {complementary arc: $Y_J$};
\end{tikzpicture}

\caption{Selected equality links in one fixed cyclic word with eight diagonal
blocks and $\ell=3$. The short arc $0\to1\to2\to3$ and its complement
indicate the positions contributing to $X_J$ and $Y_J$. The two marks in
block $5$ are distinct mate occurrences. Further occurrences and equality
links are omitted; distinct prescribed classes may use equal numerical
partition indices after inequalities between classes have been removed.}
\label{fig:shortest-pair}
\end{figure}

State Cauchy--Schwarz gives
\begin{equation}
 \left|\sum_J\varphi(X_JY_J)\right|
 \le
 \left(\sum_J\varphi(X_JX_J^*)\right)^{1/2}
 \left(\sum_J\varphi(Y_J^*Y_J)\right)^{1/2}.
 \label{eq:shortest-pair-cs}
\end{equation}
The second square sum is uniformly bounded. First drop $E_\ell'$ from
the positive operator bound for $Y_J^*Y_J$. Order the exterior mate
occurrences from left to right, using commutativity for those in one
diagonal block. Write $\mathbf j=(j_1,\ldots,j_{\ell-1})$. At the first
remaining mate occurrence, factor
\[
 B_{\mathbf j}=C_0F_jB'_{\mathbf j'},
\]
where $j$ is that occurrence's index, $\mathbf j'$ lists the others,
$\{F_j\}_j$ is its projection partition, and $C_0,B'_{\mathbf j'}$ are
independent of $j$. Then
\[
 \sum_j B_{\mathbf j}^*B_{\mathbf j}
 = {B'_{\mathbf j'}}^*
       \left(\sum_jF_jC_0^*C_0F_j\right)B'_{\mathbf j'}
 \leq\|C_0\|^2{B'_{\mathbf j'}}^*B'_{\mathbf j'}.
\]
Iterating over the remaining mate occurrences, with the fixed norm bounds
on the intervening multipliers, gives
\[
 \sum_{j_1,\ldots,j_{\ell-1}}
 B_{j_1,\ldots,j_{\ell-1}}^*
 B_{j_1,\ldots,j_{\ell-1}}
 \leq C\,1
\]
for a fixed constant $C$. When $\ell=1$, this is simply the fixed norm
bound for the exterior word. Finally sum the $E_0$ corners in the state.
As $\{\lambda_0(e_{j_0})\}_{j_0}$ is a partition in the centralizer, this
leaves a bound independent of $\mathcal P$ and of all auxiliary unitaries.

We have reduced the moment estimate to the first square sum in
Equation~\eqref{eq:shortest-pair-cs}. Its smallness will come from a
collision between the two endpoint indices. The interior pinchings
first converge at fixed coarse endpoint data; only afterward will
those coarse endpoint data be refined.

Remove $V_0$ from the state by centrality and sum the interior indices
in $X_JX_J^*$. The result is
\begin{equation}
 \sum_{j_0}
 \varphi\!\left(
   \lambda_0(e_{j_0})
   C_{\mathcal P}\bigl(\lambda_\ell(e_{j_0})\bigr)
   \lambda_0(e_{j_0})\right),
 \label{eq:endpoint-collision}
\end{equation}
where the positive map $C_{\mathcal P}:A\to M$ is given by the nested
operations
\begin{align*}
 C_{\mathcal P}(f)&=L_1,\\
 L_\ell&=y_\ell f y_\ell^*,\\
 L_h&=y_h S_{\lambda_h(\mathcal P)}
                (V_hL_{h+1}V_h^*)y_h^*
       \qquad(h=\ell-1,\ldots,1).
\end{align*}
For $\ell=1$ this means simply $C_{\mathcal P}(f)=y_1fy_1^*$.

Fix a coarse partition $\mathcal Q$ from the generating sequence and
let $\mathcal P$ refine it. Positivity permits each input projection
in Equation~\eqref{eq:endpoint-collision} to be enlarged to its coarse
piece. Regrouping the output corners by centrality then bounds that
expression by
\begin{equation}
 \sum_{e\in\mathcal Q}
 \varphi\!\left(
    \lambda_0(e)C_{\mathcal P}\bigl(\lambda_\ell(e)\bigr)
    \lambda_0(e)\right).
 \label{eq:coarse-endpoint-bound}
\end{equation}
For each fixed $f\in A$, the nested expression defining
$C_{\mathcal P}(f)$ converges in state norm to the expression in which
each pinching is replaced by $E_A$ and the $V_h$ disappear. This
convergence is uniform over the auxiliary unitaries. Indeed, work from
the innermost pinching outward. Each translated sequence
$\lambda_h(\mathcal P_j)$ is generating, so
Lemma~\ref{lem:pinching} gives its state-norm convergence to $E_A$ on
a fixed input. Conjugation by $V_h$ commutes with the pinching,
preserves the state norm, and fixes the limiting $A$-valued datum.
Pinching is contractive in the state norm. Its error can thus be
propagated through the fixed left and right multipliers $y_h,y_h^*$
using Lemma~\ref{lem:state-multipliers}, uniformly over all the
remaining phases. This proves the asserted inside-out convergence.

Define positive kernel maps on $A$ by
\[
 K_h(f)=E_A(y_hfy_h^*)\qquad(1\le h\le\ell).
\]
For each fixed input $f=\lambda_\ell(e)$, $e\in\mathcal Q$, the preceding
convergence applies. Since $\mathcal Q$ is finite, the limit of
Equation~\eqref{eq:coarse-endpoint-bound}, with $\mathcal Q$ fixed, is
\begin{equation}
 \sum_{e\in\mathcal Q}
   \int_X\lambda_0(e)
       (K_1\circ\cdots\circ K_\ell)
          \bigl(\lambda_\ell(e)\bigr)\,d\mu.
 \label{eq:coarse-kernel-collision}
\end{equation}
Proposition~\ref{prop:atomless-kernels}, applied to the adjoints of the
$y_h$, says that these kernels and their composition are atomless
almost everywhere. In taking compositions, a preceding kernel avoids
the fixed null exceptional set where a following kernel may fail to be
atomless.

Finally refine $\mathcal Q$ along the generating sequence. The
integrand sum in Equation~\eqref{eq:coarse-kernel-collision} tests
whether $\gamma_0$ of the output site and $\gamma_\ell$ of the input
site lie in the same partition piece. These conditions decrease to
equality of those two points. Since $\gamma_\ell$ is one-to-one,
this is a graph having zero conditional mass for each atomless row
of the composed kernel. Nonsingularity makes all fixed null exceptions
irrelevant. The kernel masses are bounded, so bounded convergence shows
that Equation~\eqref{eq:coarse-kernel-collision} tends to zero.

For each fixed coarse $\mathcal Q$, the limsup of the first square sum,
even after taking the supremum over the auxiliary unitaries at each
fine $\mathcal P$, was bounded by
Equation~\eqref{eq:coarse-kernel-collision}. First take that fine-partition
limsup and then refine $\mathcal Q$. The first square sum therefore
vanishes uniformly in the auxiliary phases. The second square sum was
uniformly bounded, so Equation~\eqref{eq:shortest-pair-cs} proves
vanishing of the equality-pattern sum. Averaging the auxiliary phases,
undoing the finite inclusion--exclusion and Fourier expansions, and
then removing the fixed analytic approximations proves the required
vanishing in Equation~\eqref{eq:cyclic-centered-word}. The induction
now proves Equation~\eqref{eq:moment-transfer}.
\end{proof}

The moment-transfer statement uses a fixed finite symbol alphabet.
The following approximation allows it to be applied to the measurable
color partitions supplied by Theorem~\ref{thm:bernoulli-coloring}.

\begin{lemma}[Cylinder approximation of projection partitions]
\label{lem:cylinder-partitions}
In the Bernoulli extension of Lemma~\ref{lem:bernoulli-model}, let
$p_1,\ldots,p_r\in B$ be projections with $\sum_i p_i=1$ and
$p_ip_h=0$ for $i\ne h$. For every $\delta>0$ there is a projection
partition $b_1,\ldots,b_r\in B$, each of whose members is a finite
cylinder function for one common finite dyadic symbol size, such that
\[
 \|b_i-p_i\|_\Phi<\delta\qquad(1\le i\le r).
\]
Every evaluation of this cylinder partition by the preceding
finite-symbol substitution gives a projection partition of $1$ in
$A$, modulo the fixed null sets of its definition.
\end{lemma}

\begin{proof}
The extension $\sigma$-algebra is generated by base measurable sets and
countably many label-coordinate readings along the partial graphings.
Dyadic approximations to the label coordinates therefore show that the
finite cylinder set algebra, allowing arbitrary base measurable sets,
is dense in probability. Approximate each of the measurable sets
corresponding to the $p_i$ by sets in this algebra. Disjointize the
approximations in order and put the remaining complement into one
piece, retaining exactly $r$ pieces and covering the identity. The
resulting symmetric-difference measures can be made arbitrarily small,
which gives the displayed state-norm bounds for their projections.
There are only finitely many pieces and coordinates, so all of them
use one common finite dyadic symbol size after refinement.

The resulting identities $b_i=b_i^*=b_i^2$, $b_ib_h=0$ for $i\ne h$,
and $\sum_i b_i=1$ are identities of finite cylinder functions. The
full-support argument in the construction of $\mathrm{ev}_{\mathcal P}$
shows that substitution preserves all these identities. Thus their
evaluations form a genuine projection partition in $A$.
\end{proof}
